\documentclass[12pt,a4paper]{ctexart}

\usepackage{geometry}
\geometry{left=2.5cm,right=2.5cm,top=2cm,bottom=2cm}
\usepackage{xcolor}
\newcommand{\tb}{\textcolor{red}{\textbf{【待补充】}}}

\setlength{\parindent}{0pt}
\setlength{\parskip}{0.8em}

\begin{document}

\begin{center}
{\Large \textbf{讲解讲稿 · 多种算法求解0/1背包问题的性能比较}}

\vspace{0.3cm}
讲解时长：5--10 分钟 \quad 幻灯片页数：13
\end{center}

\vspace{1cm}

% ====== Slide 1 ======
\section*{Slide 1 — 封面}

各位老师、同学大家好。我们小组的选题是2.3问题求解类，具体问题是0/1背包问题。我们的实验实现了四种算法策略，在三个数据源上进行了系统性的对比分析。接下来由我来汇报实验的设计思路和主要发现。

% ====== Slide 2 ======
\section*{Slide 2 — 0/1背包问题}

首先简单回顾一下0/1背包问题。给定n个物品，每个物品有自己的重量和价值，背包有一个固定的容量上限W。每件物品要么整件拿走，要么不拿，不能切割。目标是选出总重量不超过W的一个子集，使得总价值最大。这个问题是NP-hard的，但具有最优子结构性质，因此在算法教学中是一个经典的对照实验平台，在工业上也有资源分配、投资决策等实际应用。

% ====== Slide 3 ======
\section*{Slide 3 — 四种求解策略}

我们选择了四种算法策略来求解同一个问题。贪心算法按性价比排序，从高往低抢。动态规划用自底向上的状态转移填表，始终拿到最优解。回溯算法对解空间树做深度优先搜索，利用上界剪枝来减少搜索量。分支限界法则用优先队列来控制搜索顺序，配合上下界做更激进的剪枝。四种算法的基本理念完全不同，这为后续的对比分析提供了丰富的讨论空间。

% ====== Slide 4 ======
\section*{Slide 4 — 贪心算法}

贪心算法是这个实验里最简单的策略。每件物品算一个性价比，也就是价值除以重量，然后从高到低排序，按顺序尝试放入——只要当前容量还够就塞进去，不够就跳过。时间复杂度是n log n，排序主导，实际操作与容量W完全无关，因此无论W多大都能秒出结果。但代价是它不保证最优解。失效的场景非常直观：一件高性价比但很重的东西塞进去之后，剩余空间装不进其他物品了，而本来可以选两件稍低性价比但加起来更值钱的组合。贪心不会"后悔"，所以这个gap在某些数据上可以到15\%。

% ====== Slide 5 ======
\section*{Slide 5 — 动态规划}

动态规划是这四种算法中唯一既保证最优解、又在大多数数据集上跑得动的方法。状态定义是dp[w]等于容量为w时的最大价值。转移方程就是著名的max(dp[w], dp[w-wᵢ]+vᵢ)。实现上我们使用了一维滚动数组，把空间降到O(W)。这里有一个很容易出错的细节——内层容量循环必须从大到小倒序遍历，否则正序会导致同一物品被多次使用，变成了完全背包问题，结果就是错的。DP的问题是复杂度O(nW)是伪多项式——表面看n乘W是多项式，但W是数值不是输入长度，W翻10倍运行时间就翻10倍。我们的实验里W等于250万的时候DP大概跑10秒，W到1亿的时候跑了将近两分钟，W到100亿的时候一个int数组就40个G，直接内存溢出。

% ====== Slide 6 ======
\section*{Slide 6 — 回溯与分支限界}

回溯与分支限界部分由张诗淇负责实现。回溯在做深度优先搜索时通过当前重量和价值计算剩余容量的上界，如果上界还不如已知最优解就直接剪掉整个子树。分支限界法在这个基础上用优先队列取代递归栈，优先扩展上界最高的节点，理论上有更好的剪枝效率。具体的策略描述、复杂度分析和实验结果请张诗淇补充。

% ====== Slide 7 ======
\section*{Slide 7 — 三个数据源}

我们的实验使用了三个数据源，从不同维度覆盖问题空间。第一个是FSU，Florida State University提供的8个基准实例，n从5到24，规模虽小但有官方最优解，适合用来验证算法正确性。第二个是我们参考Pisinger的方法自生成的Pi数据集，3种重量-价值相关性模式乘以3种规模，共9个实例，用于控制变量地分析数据分布和规模各自的独立影响。第三个是JorikJooken在GitHub上发布的学术标准测试集，我们选取了其中6个实例，n从400到1000，W从100万到100亿，用于测试算法在接近真实场景下的性能表现。

% ====== Slide 8 ======
\section*{Slide 8 — FSU 实验结果}

先看FSU的结果。DP全部8个实例命中官方最优解，证明了我们代码的正确性。贪心这边就比较惨了，8个里面只有P01这一个达到了最优，其余都有偏离。最夸张的是P06，贪心差了14.81\%——7个物品的简单问题，贪心就因为拿了一件高价值的东西导致后面好的组合装不下。而P08有24个物品、容量640万，DP跑了1.59秒，贪心0.12毫秒，速差超过一万倍。这个对比直观地展示了两种算法策略的代价差异。

% ====== Slide 9 ======
\section*{Slide 9 — Pi 实验结果}

Pi实验是我们最有系统性的分析。关键发现是：贪心的表现完全取决于数据分布。不相关数据——就是重量和价值各自独立随机——此时物品之间的性价比天差地别，排序以后好坏一目了然，贪心gap几乎为零。而到了强相关数据——所有物品的价值约等于重量加一个常数，性价比差异极小——此时排序基本没有区分度，贪心在50个物品的小数据上就差了0.85\%。DP这边，耗时从50个物品的10毫秒暴涨到2000个物品加250万容量的将近10秒，和n乘W的增长严格一致，验证了伪多项式的理论预测。贪心呢？始终在毫秒级。

% ====== Slide 10 ======
\section*{Slide 10 — JJ 实验结果}

JJ实验最有意思的不是c等于100万那6个实例的数据——它们和Pi的结论完全一致，贪心gap都在0.05\%以内——而是在于极限测试。我们测了一个c等于1亿的实例，DP跑了113秒，而贪心不到1毫秒，差距11万倍。c等于100亿的时候DP直接OOM，因为一个int数组就40G，普通电脑根本分配不出来。回溯在n等于400的时候搜索空间已经是2的400次方量级，即使有上界剪枝也完全跑不动。这些极限数据远比理论分析有说服力——它们证明了各种算法的瓶颈不是理论推导出来的假设，而是实验台上真实发生的。

% ====== Slide 11 ======
\section*{Slide 11 — 核心洞察}

这一页是我们整个实验最核心的结论。把四种算法放进一个二维坐标系里看——横轴是n的大小，纵轴是W的大小——就能看清楚各自的领地。左下角n小W小，谁跑都行。右下角n大W小，DP是王者，因为它的复杂度只乘W，W小就快；回溯在n大的时候指数爆炸，完全不可行。左上角n小W大，反过来回溯和分支限界有优势，因为n小搜索空间可控，而DP被大W拖慢甚至OOM。右上角n大W大，对不起，只有贪心能跑——虽然它不保证最优解，但至少能给你一个答案。所以结论不是"哪个算法好"，而是"在什么条件下该用哪个"。算法选型从来都是工程权衡。

% ====== Slide 12 ======
\section*{Slide 12 — 复杂度对比}

这张表是对前面所有讨论的形式化总结。贪心是n log n排序、O(n)空间、不保证最优、怕强相关数据。DP是n乘W伪多项式、W空间、保证最优、怕大容量。回溯和分支限界这边请张诗淇补充具体的复杂度分析和敏感性特征。但从我们已经完成的部分已经可以看出一个清晰的模式：算法设计就是一个在时间、空间和正确性三者之间的三维权衡，没有免费午餐。

% ====== Slide 13 ======
\section*{Slide 13 — 总结}

最后总结五个要点。第一，0/1背包问题是检验算法策略优劣的理想平台，因为它的结构足够简单所以各种算法都能往上套，但又足够难所以能拉开差距。第二，贪心毫秒级别，与W无关，但不保证最优，gap可以到15\%。第三，DP始终最优，但O(nW)在大容量下不可用，100亿直接OOM。第四，回溯和分支限界介于二者之间，n小时有可能优于DP。第五，选型原则就是看n和W的相对大小，以及应用场景对精度的需求。我们想用一句话收尾：时间与最优性的折中，是算法设计领域永恒的主题。谢谢大家。

\end{document}
